\centerline{\bf \Huge ДЗ$\cdot$Матлогика}
\title{ДЗ$\cdot$Матлогика}
\begin{flushright}
    \scriptsize $-$Сейчас ты скажешь «Эльзята, как ты узнала про фальшивую монету?!»\\
    $-$Эльзята, как ты узнала про фальшивую монету?!\\
    $-$Yare Yare Daze\\
    Имя стенда $-$ Фальшивомонетчик. Имя мастера $-$ НеНачальник.\\
\end{flushright}
\problem{1} Расул Владимирович решил наградить Эльзяту за успешное выполнение домашнего задания и отправил ей 33 одинаковые монеты. По пути эти монеты попали в руки фальшивомонетчика НеНачальника, который заменил одну из них на похожую, но меньшую по весу. За какое минимальное количество взвешиваний на чашечных весах Эльзята определит, какая монета фальшивая? Приведите полный алгоритм действий Эльзяты!

\problem{2}Как-то на встрече ЛМШ Яне, Алине, Герману и Саше нужно было выстроиться в шеренгу по старшинству, начиная с младшего. Начальник ЛМШ не знал точный возраст ребят, и они решили над ним подшутить.
Каждый из них сказал:\\
Герман: «Яна младше всех, а я старше»,\\
Яна: «Алина старше Саши и Германа»,\\
Алина: «Яна старше Германа, Саша врет»,\\
Саша: «Я старшая среди девочек».\\
Начальник точно знает, что двое сказали правду, а другие двое соврали, и ему не составило особого труда определить все варианты расстановок. А сможете ли вы?

\problem{3} В компанию ЛмШ поступили новые работники: Шиншиллы и Лососи. Так как Шиншиллы еще не отучились, а Лососи еще не научились лгать. Шиншиллы говорят только ложь, а Лососи - только правду.
В первый рабочий день надо было рассортировать новых сотрудников, так как вся компания была занята подготовкой к турниру, свободен был только старый Лосось, которого подводило зрение. И чтобы понять, сколько среди новых 28 сотрудников Лососей, а сколько Шиншилл, он построил всех в ряд и поочереди задавал два вопроса: среди вас Шиншилл больше 14, среди вас Лососей меньше 16?(Первому - среди вас Шиншилл больше 14, второму - среди вас Лососей меньше 16, третьему - среди вас Шиншилл больше 14 и так далее…)
Оказалось, что на первый вопрос 8 сотрудников ответили - да (соответственно 6 ответили - нет), а на второй вопрос 9 сотрудников ответили - да (соответственно 5 ответили - нет).
Помогите старому Лососю узнать, сколько Шиншилл и сколько Лососей, среди новеньких сотрудников.